% =====================================================================
% Figura: planteamiento del problema de contorno (recinto V, frontera S).
% Estilo de los apuntes de Mecánica: sin marco, superficies azules
% translúcidas y trazos finos. Fuentes rayadas: conocidas dentro de V
% (celeste, contorno continuo) y desconocidas fuera (gris, contorno
% punteado). Vectores r, r', R; elemento dS' con su normal exterior.
% =====================================================================
\begin{tikzpicture}[scale=1.0]
  % ---------------------------------------------------- recinto V y S
  \path[fill=azul,fill opacity=0.17,draw=azul,line width=0.75pt]
    plot[smooth cycle,tension=0.75] coordinates
    {(0,0) (2.2,-0.5) (4.6,0.2) (5.2,2.0) (3.8,3.3) (1.4,3.0) (-0.3,1.8)};
  % ---------------------------------------------------- fuentes conocidas (en V)
  \tikzset{fuenteV/.style={pattern={Lines[angle=45,distance=2.3pt,line width=0.45pt]},
      pattern color=celeste,draw=celeste,line width=0.55pt},
    fuenteExt/.style={pattern={Lines[angle=-45,distance=2.6pt,line width=0.4pt]},
      pattern color=gristexto!80,draw=gristexto!80,line width=0.5pt,
      dash pattern=on 1.6pt off 1.2pt}}
  \path[fuenteV] plot[smooth cycle,tension=0.8] coordinates
    {(0.72,1.15) (1.05,0.82) (1.6,0.92) (1.72,1.38) (1.3,1.66) (0.85,1.55)};
  \path[fuenteV] plot[smooth cycle,tension=0.8] coordinates
    {(3.62,0.72) (3.95,0.5) (4.3,0.7) (4.22,1.02) (3.85,1.08)};
  \node[rotulo,text=celeste,anchor=east] at (0.62,1.25) {$\rho(\vrp)$};
  % ---------------------------------------------------- fuentes exteriores
  \path[fuenteExt] plot[smooth cycle,tension=0.8] coordinates
    {(6.6,3.05) (7.0,2.75) (7.55,2.95) (7.45,3.4) (6.95,3.5)};
  \path[fuenteExt] plot[smooth cycle,tension=0.8] coordinates
    {(6.35,-0.55) (6.7,-0.8) (7.05,-0.55) (6.85,-0.25) (6.5,-0.25)};
  \path[fuenteExt] plot[smooth cycle,tension=0.8] coordinates
    {(-1.25,3.25) (-0.95,3.0) (-0.55,3.2) (-0.7,3.5) (-1.1,3.55)};
  \node[rotulogris,anchor=north] at (6.75,-0.95) {fuentes exteriores};
  % ---------------------------------------------------- elemento dS' y normales
  \draw[celeste,line width=1.3pt,line cap=butt]
    ($(5.2,2.0)+(-0.25*0.27,-0.25*0.963)$) -- ($(5.2,2.0)+(0.25*0.27,0.25*0.963)$);
  \draw[vectorgris,draw=celeste] (5.2,2.0) -- ++(0.963*0.75,-0.27*0.75)
    node[right,rotulo,text=celeste] {$\un$};
  \node[rotulo,text=celeste,anchor=west] at (5.32,2.38) {$\dd S'$};
  \draw[vectorgris,draw=celeste] (1.4,3.0) -- ++(-0.17*0.7,0.985*0.7)
    node[above,rotulo,text=celeste] {$\un$};
  % ---------------------------------------------------- vectores de posición
  \coordinate (O) at (-0.9,-1.55);
  \coordinate (rp) at (1.2,1.25);
  \coordinate (r) at (3.9,2.15);
  \draw[gristexto,line width=0.5pt,-{Stealth[length=4.5pt,width=3.6pt]}] (O) -- (rp)
    node[pos=0.45,left=1pt,rotulo,text=gristexto] {$\vrp$};
  \draw[gristexto,line width=0.5pt,-{Stealth[length=4.5pt,width=3.6pt]}] (O) -- (r)
    node[pos=0.42,below right=-1pt,rotulo,text=gristexto] {$\vr$};
  \draw[vector] (rp) -- (r) node[pos=0.47,above=1pt,rotulo,text=azul] {$\vR=\vr-\vrp$};
  \fill[texto] (O) circle (1.3pt) node[left=1pt,rotulogris] {$O$};
  \fill[texto] (rp) circle (1.3pt);
  \fill[texto] (r) circle (1.6pt) node[right=2pt,rotulo] {$P$};
  % ---------------------------------------------------- rótulos del recinto
  \node[rotulo,text=azul] at (2.6,2.7) {$V$};
  \node[rotulo,text=azul,anchor=west] at (4.75,0.05) {$S=\partial V$};
\end{tikzpicture}
